\section{Introduction}
Multi-task learning optimizes several objectives with shared parameters \cite{caruana1997multitask}. These objectives often differ in numerical scale, curvature, and noise, as in dense prediction with segmentation, depth, and surface-normal losses or tabular problems with multiple binary and regression targets. A weighting rule that responds directly to raw loss magnitude can let one task dominate the shared update and push adaptive uncertainty variables toward uninformative regimes.

Classical homoscedastic uncertainty weighting learns one log-variance per task \cite{kendall2018}. Although effective in many settings, its unconstrained coordinates can turn severe loss-scale mismatch into large precision shifts. Methods such as PCGrad \cite{yu2020pcgrad} address gradient conflict instead, leaving a gap for adaptive weighting designed specifically for heterogeneous task scales.

We study Bounded Precision-Geometry Scaling (BPGS), which maps one learnable uncertainty coordinate per task through a bounded sigmoid chart and, in its canonical form, anchors the resulting log-variance to detached batch log-loss statistics. Network parameters use detached normalized weights, while uncertainty coordinates are updated with a separate objective. Our question is narrow: does this bounded, batch-aware parameterization improve robustness to scale mismatch while remaining competitive on standard benchmarks?

Our contributions are:
\begin{enumerate}[leftmargin=1.25em]
\item We formulate BPGS as a bounded, batch-aware, split uncertainty-weighting method and prove normalized-weight invariance to uniform loss rescaling when safeguards are inactive.
\item On pure loss rescaling, BPGS remains effectively invariant from $\times 1$ to $\times 1000$ (macro score 0.777 to 0.778), whereas Kendall drops from 0.780 to 0.637; L1 normalization does not recover this stability.
\item On NYUv2, BPGS achieves the best depth metrics and lowest total loss among the compared methods, including Nash-MTL; ablations isolate the batch-aware chart, first-batch calibration, and split stop-gradient.
\item Controlled NYUv2 diagnostics measure calibration and batch-size sensitivity plus runtime and memory overhead, while Yeast and RF1 test transfer beyond synthetic stress settings.
\end{enumerate}
